Task 3 được trình bày theo diễn tiến thực tế của dự án: yêu cầu phần mềm, ba increment, tích hợp và triển khai, sau đó là các hoạt động quản lý xuyên suốt. Bảng \ref{tab:t3-map} chỉ ra nơi đáp ứng từng yêu cầu của đề bài.

\begin{table}[H]
\centering\small
\caption{Đối chiếu yêu cầu Task 3 với các mục của báo cáo}\label{tab:t3-map}
\begin{tabularx}{\linewidth}{@{}p{.30\linewidth}X@{}}
\toprule
\textbf{Yêu cầu của đề bài} & \textbf{Mục trong báo cáo}\\\midrule
3.1 Triển khai phạm vi phần mềm & Yêu cầu phần mềm (Mục \ref{sec:t3-requirements}), mô hình use case (Mục \ref{sec:t3-usecase}), Increment 1--3 (Mục \ref{sec:i1-development}, \ref{sec:i2}, \ref{sec:i3}), tích hợp và triển khai (Mục \ref{sec:t3-deploy}).\\
3.2 Nhật ký dự án hàng tuần & Mục \ref{sec:t3-weekly}.\\
3.3 Theo dõi so với Baseline V1.0 & Mục \ref{sec:t3-actual} và phần ``Tiến độ và checkpoint'' của mỗi increment.\\
3.4 Cập nhật sổ rủi ro & Mục \ref{sec:t3-risk} và phần ``Vấn đề phát sinh'' của mỗi increment.\\
3.5 Kiểm soát chất lượng & Mục \ref{sec:t3-quality} và phần ``Kiểm thử'' của mỗi increment.\\
3.6 Phản hồi thay đổi của giảng viên & Mục \ref{sec:t3-change}.\\
3.7 Project Baseline V2.0 & Mục \ref{sec:t3-baseline-v2}.\\
3.8 Kế hoạch chi tiết & Phần ``Kế hoạch chi tiết'' của mỗi increment và của tích hợp.\\
\bottomrule
\end{tabularx}
\end{table}
